\begin{proof}[The refined class of the weighted OS45 member]
\label{maps:proof:D.OS.refined}
We retain the exact source sections, fan, line linearization and peripheral
markings in the statement.  In particular, no translation of a selected fan
edge, and no replacement of a period unit by its valuation, is made below.
Write $C$ for the integral sum of the actual complex critical curves.
\begin{enumerate}
\item \emph{The source line and the selected cusp chart.}
Let $q_E$ be the based multiplicative elliptic period and use the actual
multiplicative lifts of $P$ and $Q$.  The calculation underlying
\eqref{maps:eq:D.FI2} gives
\[
 \sigma(Z)=\frac{\theta(Z/P;q_E)}{\theta(Z;q_E)},\qquad
 \sigma(q_E Z)=P\sigma(Z),\qquad e_{\rm flat}=s_{P-O}/\sigma .
 \tag{D.OR1}\label{maps:eq:D.OR1}
\]
Here $s_{P-O}$ is the canonical meromorphic section of
$\mathcal O(P-O)$.  If $W$ is its nonzero-line coordinate in $e_{\rm flat}$,
the manuscript coordinate is $V=W^{-1}$.  This is the tautological map from
the nonzero line to its dual; it inverts every line multiplier.

At the $I_8$ site the actual lifts have valuations $6$ and $24$.
The extending-frame heights are $F_i=\max(6-i,0)$ for $0\leq i<8$,
with $F_{i+8}=F_i-6$.  Thus the source-line graph has the edge
$(0,6)$--$(1,5)$.  It is an edge of the prescribed triangular fan.
The adjacent ordered cone
\[
 (1,0,5),\quad(1,1,5),\quad(1,0,6)
\]
has determinant one.  Its literal coordinates, and the actual extending
physical line coordinate, are
\[
 q=xyz,\quad Z=y,\quad V=x^5y^5z^6,\quad
 W_{\rm ext}=\frac{q^6}{ZV}=x,\qquad
 (X,Y,W_{\rm ext})=(xz,y,x).
 \tag{D.OR2}\label{maps:eq:D.OR2}
\]
The inverse is $(x,y,z)=(W_{\rm ext},Y,X/W_{\rm ext})$ and its forward
Jacobian is $-x$.  In particular the line winding has not been discarded.
The extending frame is exactly $(Z/q^6)e_{\rm flat}$: its orders on both
branches vanish, $P$ and $O$ miss the node, and the frame and its inverse
extend over the remaining codimension-two point.

The marked weights in this cone are
\[
 e_1=(-1,1,0),\qquad \delta=(-1,0,1),\qquad
 a_0=e_1-\delta=(0,1,-1).
 \tag{D.OR3}\label{maps:eq:D.OR3}
\]
Thus $\delta$ is inverse physical line rotation and $a_0$ is the source
vanishing circle at constant $W_{\rm ext}$.  At the first $I_1$ site its
zero section cocycle gives the extending flat frame and the vanishing
circle $a_1=e_1$.  Both oriented boundary areas are
$a_i\wedge\delta=\alpha=e_1\wedge\delta$; changing the common boundary
orientation changes both signs.  The two meridians $(1,-1),(1,0)$ have
determinant one.

The cusp period lattice has columns $(8,-6),(24,-16)$ and is
$8\mathbb Z\oplus2\mathbb Z$.  Translating the selected edge by $(0,-6)$
would use $9(8,-6)-3(24,-16)$, and hence the actual source multiplier
$q_E^9Q^{-3}$ and its actual line multiplier.  No such translation is used.

\item \emph{The based collar and its embedded odd sphere.}
Choose a narrow neighborhood of the prescribed based arc from this cusp
to the first $I_1$ site, containing no other singular site.  The two based
elliptic matrices are $T^8,T$ and their $G$ cocycles are $(3,0),(0,0)$.
They fix the same short circle.  Thus the based exponentials $q_E,P,Q$
are single-valued on this punctured neighborhood, with valuations
$(8,6,24)$ and $(1,0,0)$ at the two ends.  If the actual forward suspension
lift has coefficient $c_F(Z)$ in the meromorphic section, its coefficient
in the flat frame is
\[
 \kappa(t)=c_F(Z)\frac{\sigma(QZ)}{\sigma(Z)}.
 \tag{D.OR4}\label{maps:eq:D.OR4}
\]
The cancellation of its divisor and automorphy makes this a nonzero
holomorphic function on the compact elliptic fibre, hence independent of
$Z$.  Its valuations are $16,0$, with the inverse-lift convention of
\eqref{maps:eq:D.FI2}: $w=-16$ and $8w-(-6)24=16$ at the cusp.
The actual quotient maps are therefore
\[
 (Z,V)\longmapsto(q_E Z,P^{-1}V),\qquad
 (Z,V)\longmapsto(QZ,\kappa^{-1}V).
 \tag{D.OR5}\label{maps:eq:D.OR5}
\]
All their phases and nonzero holomorphic units are retained.  The origin
is the received section $O$, namely $Z=1$, and the frame is
\eqref{maps:eq:D.OR1}; these are also the endpoint choices.

Parametrize the embedded based arc by a smooth regular $p:[0,1]\to B$,
linear in affine coordinates near its ends.  On a sufficiently thin strip
the explicit map $t(s+iu)=p(s)+iu p'(s)$ is an orientation-preserving
tubular diffeomorphism; along the real interval its Jacobian is
$|p'(s)|^2$.  Retain the actual nonzero endpoint units in
$q_0(t(f))=fc_0(f)$ and $q_1(t(f))=(1-f)c_1(f)$.
On $\mathcal Q=\{ad-bc=1\}$, put $f=ad$ and define
\[
 \Phi(a,b,c,d)=[\,t(f),\ Z=ab,\ V=f^6/a\,],
 \tag{D.OR6}\label{maps:eq:D.OR6}
\]
where the brackets mean exactly \eqref{maps:eq:D.OR5}.  At the cusp this
extends in the selected fan chart by
\[
 (x,y,z)=(c_0^6/b,\ ab,\ c_0^{-5}d);
\]
at the other endpoint it extends in the original source-node/line chart by
\[
 (X,Y,W)=(-c_1c/a,\ ab,\ a/f^6).
 \tag{D.OR7}\label{maps:eq:D.OR7}
\]
The core Jacobians are respectively $c_0/b$ and $c_1$, hence nonzero.
On the middle collar the fibre character matrix on $(a,b)$ is
$\left(\begin{smallmatrix}1&1\\-1&0\end{smallmatrix}\right)$, of determinant
one, and the base derivative preserves orientation.

The compact real form
\[
 N_0=\left\{
 \begin{pmatrix}\alpha&\beta\\-\overline\beta&\overline\alpha\end{pmatrix}:
 |\alpha|^2+|\beta|^2=1\right\}
 \tag{D.OR8}\label{maps:eq:D.OR8}
\]
has generic short two-torus fibre, with the two meridians in
\eqref{maps:eq:D.OR3}.  The physical line transition is
$W_1=(Z/f^6)W_0$; its $Z$ factor has winding one and its $f^6$ factor
is retained in \eqref{maps:eq:D.OR6}.  These are the two Hopf charts of
the displayed $S^3$.
The map is injective on $N_0$: on a regular fibre the real short plane
intersects the full period lattice in exactly
$\mathbb Ze_1+\mathbb Z\delta$, by real independence of the four marked
periods; distinct arc parameters have distinct base images.  At either
end no nonzero period translation fixes the selected finite edge or core
circle.  The endpoint Jacobians prove immersion.  Compactness then gives
an embedded sphere and a neighborhood on which $\Phi$ is an embedding.

Left multiplication by
$L_\varepsilon=\left(\begin{smallmatrix}1&\varepsilon\\0&1\end{smallmatrix}\right)$,
for sufficiently small positive $\varepsilon$, gives
\[
 a=\alpha-\varepsilon\overline\beta,\quad
 b=\beta+\varepsilon\overline\alpha,\quad
 c=-\overline\beta,\quad d=\overline\alpha.
 \tag{D.OR9}\label{maps:eq:D.OR9}
\]
It misses $a=d=0$ and $b=c=0$, and the neighborhood meets no other
critical curve.  It commutes with the scalar right action, so its image
$N$ is a free scalar-invariant embedded sphere disjoint from $C$.
The unperturbed sphere has Mayer--Vietoris boundary $\alpha$.
Here is the integral detector that proves the needed oddness.
The actual matrix $A_0$ of \eqref{maps:eq:D.PC5} is the identity modulo
$2$. In the ordered bivector basis
$(e_1\wedge e_2,e_1\wedge\delta,e_1\wedge\gamma,
 e_2\wedge\delta,e_2\wedge\gamma,\delta\wedge\gamma)$, the row
$(0,1,0,0,0,0)$ modulo $2$ annihilates every column of
$\bigwedge^2A_0-I$ and takes $\alpha$ to $1$.
The Wang sequence is
\[
0\longrightarrow\operatorname{coker}(\bigwedge^2A_0-I)
\longrightarrow H_2(\partial U;\mathbb Z)
\longrightarrow\ker(A_0-I\mid H_1(T^4;\mathbb Z))
\longrightarrow0.
\]
Its right group is free, so the sequence splits as a sequence of
abelian groups. The displayed mod-two character therefore extends
to $H_2(\partial U;\mathbb Z)$. In particular $\alpha$ is not twice
any boundary class. Applying the character to the Mayer--Vietoris
image proves $[N]\notin2H_3(X;\mathbb Z)$. Since this last group is
$\mathbb Z/8$, the sphere is an odd generator. The perturbation
preserves its class.

\item \emph{The exact reference frame and its local degree.}
Let $\rho$ be the reference from the abstract kernel and target bundles
in \eqref{maps:eq:D44.5}.  Fix the nowhere-zero abstract section
$\delta$.  Comparing a global second kernel vector with the local one
changes the quotient complex-line frame by a map to $\mathbb C^*$;
the choices of its lift form an affine space.  The target-frame change
also maps to $\mathbb C^*$.  On $N\simeq S^3$ these changes have zero
degree-three homotopy class.  Splittings of the differential sequence
form an affine space as well.  Thus the actual global reference and the
local reference below have zero integral degree-three discrepancy.

This comparison remains valid along the smooth middle base map.
In vertical/horizontal splittings the derivative of \eqref{maps:eq:D.OR6}
has complex-linear vertical block and real base block
$a+b\,\text{conjugation}$ with $|a|>|b|$.  Replacing the latter by
$a+(1-r)b\,\text{conjugation}$ and contracting the horizontal-to-vertical
block gives invertible oriented bundle maps throughout, fixed at the
holomorphic ends and retaining $\delta$ and the vertical kernel.

In the left complex Lie frame of $SL_2(\mathbb C)$ the $\delta$ vector is
constant diagonal and, on the other two coordinates, $df=ac\,u+bd\,v$.
Set $A=ac$, $B=bd$.  The Bott matrix with first column
$(B,-A)/\sqrt{|A|^2+|B|^2}$ gives the local reference, after the affine
splitting changes just described.  On \eqref{maps:eq:D.OR9},
\[
 A=-\alpha\overline\beta+\varepsilon\overline\beta^2,
 \qquad B=\beta\overline\alpha+\varepsilon\overline\alpha^2.
 \tag{D.OR10}\label{maps:eq:D.OR10}
\]
The pair has no common zero for $\varepsilon>0$: off the axes the two
equations imply $z_1=-\varepsilon^2z_1$, with
$z_1=\overline\alpha,z_2=\overline\beta$, and on each axis the remaining
term is nonzero.  Increasing $\varepsilon$ and normalizing gives the
pair $(z_1^2,-z_2^2)$.  The proper holomorphic quadratic map has four
regular preimages, all of positive real local degree.  Radial
normalization and simultaneous conjugation therefore give degree four
for this Bott matrix.  Consequently $\rho$ induces the odd framed
parity of a Lie-parallelized $S^3$.

The normalization of this parity is the quaternionic left-multiplication
frame: its stable difference from the bounding frame is one unit.
An oriented rank-three normal-frame change stabilizes with factor two.
Indeed the adjoint complex roots are $2x,0,-2x$, so
$c_2(\operatorname{ad}_{\mathbb C})=4c_2(\mathrm{fund})$ and
$\lambda(\operatorname{ad}_{\mathbb R})=-2c_2(\mathrm{fund})$;
the quaternionic rank-four representation instead has
$\lambda=-c_2(\mathrm{fund})$.
These are the precise maps in
\cite[Remark~2.17, p.~12; \S6.1, pp.~24--25]{Redden2011StringStructures}.
Thus the resulting parity is independent of the oriented normal frame.

\item \emph{The integral equivariant characteristic class.}
Put $u=c_1(\mathcal O(1))\in H^2(BS^1;\mathbb Z)$.  On each of the
sixteen fixed curves let $x$ be its positive generator.  The actual
tangent and normal data give roots
$2x,-x+u,-x-u$, with $x^2=0$.
The ordinary first Chern class is zero, and the integral Gysin sequence
of $X\times ES^1\to X_{S^1}$ makes the kernel in degree two the
multiples of $u$.  The sum of the fixed weights is zero; hence
$c_{1,S^1}(TX)=0$ and $\lambda_{S^1}=-c_{2,S^1}$.
Its restriction to every fixed curve is $u^2$.

As the ordinary $\lambda$ vanishes, that same sequence gives
\[
 \lambda_{S^1}(TX)-u^2=u\eta,\qquad
 0\longrightarrow\mathbb Zu\longrightarrow H^2_{S^1}(X;\mathbb Z)
 \longrightarrow\mathbb Z^{15}\longrightarrow0.
 \tag{D.OR11}\label{maps:eq:D.OR11}
\]
For this Gysin sequence and its Thom-pair construction see
\cite[\S4.D, pp.~438,444]{Hatcher2002}; its total space here is
homotopy equivalent to $X$ and its Euler class is the displayed $u$.
Restriction to a fixed curve gives $u\eta=0$ and hence $\eta=0$ there.
Over the rationals the Borel spectral sequence has only even entries,
so all its differentials vanish.  Multiplication by $u$ is injective
on its associated graded and hence there is no $u$-torsion.
Localization, in the precise form
\cite[Theorem~3.5 and the circle specialization, p.~8]{AtiyahBott1984},
now gives $\eta=0$ rationally.  The integral group in
\eqref{maps:eq:D.OR11} is free, so $\eta=0$ integrally:
\[
                         \lambda_{S^1}(TX)=u^2.
 \tag{D.OR12}\label{maps:eq:D.OR12}
\]
In particular no division of a torsion Pontryagin class is used.

\item \emph{The Borel pair and the actual framing transport.}
The scalar action on the sphere coordinates is
$(\alpha,\beta)\mapsto(s^{-1}\alpha,s\beta)$.
The explicit map $v(\alpha,\beta)=(\overline\alpha,\beta)$ identifies
it with the scalar Hopf sphere.  Its orientation sign is retained; only
the parity of the ensuing unit is needed.  Form
\[
 W=(N\times D^2)/((n,z)\sim(sn,sz)),\qquad
 \partial W\longrightarrow N,\quad[n,z]\longmapsto z^{-1}n.
\]
With the diagonal convention for the Borel quotient the actual map of
pairs into $(X_{S^1},X)$ is
\[
 [n,z]\longmapsto
 [\,n,(z,\sqrt{1-|z|^2}\,v(n))\,].
 \tag{D.OR13}\label{maps:eq:D.OR13}
\]
At $|z|=1$ it is $[z^{-1}n,(1,0,0)]$, exactly the ordinary inclusion
after the displayed boundary identification.  The pulled-back tangent
bundle is the pullback of $(TX|_N)/S^1$ from $N/S^1$; on the boundary
its vector identification is the actual derivative $d(z^{-1})$.

The collapsed map
\[
 W/N\longrightarrow\mathbb{CP}^2,\qquad
 [n,z]\longmapsto[z:\sqrt{1-|z|^2}\,v(n)]
\]
is a homeomorphism.  Off $[1:0:0]$, a unit representative $(Z_0,Z')$
has inverse $[v^{-1}(Z'/\|Z'\|),Z_0]$; changing the unit representative
by a phase is exactly the defining equivalence in $W$.
The collapsed boundary supplies the remaining point.
Thus the square of the base generator evaluates to a unit on $(W,N)$.

The stable spin tangent bundle of $X$ has a global stable frame $\tau$:
in dimension six its only obstruction is $\lambda=0$.
The required stable homotopy range is precisely the orthogonal period
in \cite[Theorem~II and (1.5), pp.~314--315]{mapsBott59}, with the unit
generator fixed by \eqref{maps:eq:D44.6}.
Choose a point away from the sphere, the critical curves and the chains
under discussion.  On the punctured manifold, which has a
five-dimensional spine, rank-six cancellation realizes this stable
trivialization by an actual tangent frame.  Its comparison on $N$
has the same integer degree, since
$\pi_3\operatorname{Spin}(6)\to\pi_3\operatorname{Spin}$ is an isomorphism.
All applications to actual admissible three-frames and their complements
below take place on this punctured manifold.  No tangent framing of the
closed manifold is asserted.
The frame defines a relative lift $\lambda_\tau$ on $(X_{S^1},X)$.
By \eqref{maps:eq:D.OR12}, its difference from the canonical relative
lift of $u^2$ lies in the image of $H^3(X;\mathbb Z)=\mathbb Z/8$.
Pullback to $H^4(W,N;\mathbb Z)=\mathbb Z$ kills that ambiguity.
Therefore the global frame differs from a quotient-descending frame
on $N$ by an odd integral stable framing degree.

The quotient bundle is spin over $S^2$, hence trivial there.
The quotient of $TN$ is $\mathbb R\delta\oplus TS^2$, and its rank-three
normal quotient is also spin and trivial over $S^2$.
Thus a quotient-descending frame induces the Hopf tangent framing,
whose parity is one by the quaternionic computation in step~3.
Changing its spin gauge over $S^2$ contributes nothing, since
$\pi_2\operatorname{Spin}=0$; changing the normal frame contributes
an even stable degree by the explicit adjoint map in step~3.
The stable tangent framing induced by $\tau$ on $N$ consequently has
parity $1+1=0$.

\item \emph{The complete critical-chain intersection.}
The preceding global/local comparison also determines the required
integer intersection parity, rather than inferring it from the core
coefficients.  The exact reference in step~3 and a quotient-descending
frame have zero degree-three discrepancy: the reference Bott matrix
$A_2$ satisfies $A_2(sn)=D_2(s)A_2(n)$ in the left Lie frame, where
$D_j(s)=\operatorname{diag}(s^j,s^{-j})$.
It therefore descends, and any further spin-frame comparison is over
$S^2$ and has zero degree-three class.  Equivalently, the literal matrices
\[
 A_1=\begin{pmatrix}\overline\alpha&-\beta\\
                     \overline\beta&\alpha\end{pmatrix},\qquad
 C_1=A_2A_1^{-1},\qquad
 C_1(sn)=D_2(s)C_1(n)D_1(s)^{-1}
\]
have degrees $1,3$ respectively; this is the actual weight-two to
weight-one frame map, while the reference differs by $A_1^{-1}$.
Together with step~3 this proves that the global frame and the actual
$\rho$ differ by an odd integer on $N$.

The unit Bott residue \eqref{maps:eq:D44.6} on every actual double curve
gives the comparison class of \eqref{maps:eq:D44.7} on $X\setminus C$.
Let $D$ be the sum of the prescribed critical-specialization chains
and the complete integral core nullchain \eqref{maps:eq:D.OS8}; thus
$\partial D=C$.  Its Thom intersection cochain has those same unit
residues.  The difference extends to a class in $H^3(X;\mathbb Z)$,
which evaluates to zero on the torsion class $[N]$.
Consequently $N\cdot D$ is odd.  This evaluates the total intersection
with the prescribed chains; it assigns no unsupported individual
local crossing numbers.  Since the core coefficients are even,
the sum of the actual specialization-chain intersections is odd.

\item \emph{The refined class.}
The reference parity is odd by step~3 and the total link is odd by
step~6.  The integral relative transgression giving
\eqref{maps:eq:D.WI13} applies to this genuinely embedded torsion sphere:
its proof uses the same unit Thom residues and an admissible sphere
frame.  Hence
\[
 \operatorname{lk}(N,\operatorname{PD}P_{\rm ref})
 =-\tfrac12\ell(N,C)-\tfrac12\epsilon_\rho=0
                  \quad\hbox{in }\mathbb Q/\mathbb Z.
\]
The sphere is odd in $H_3(X)=\mathbb Z/8$ and the torsion linking
pairing is nondegenerate.  Also $2P_{\rm ref}=\lambda=0$ kills the
free part.  We obtain
\[
                              P_{\rm ref}(X)=0.
 \tag{D.OR14}\label{maps:eq:D.OR14}
\]
The even induced framing in step~5 gives the same conclusion by
Zhubr's actual admissibility criterion
\cite[Theorem~1, p.~6; \S5, pp.~8--9]{ffZhubr13}, transported in
\eqref{maps:eq:D44.4}; the separate change-of-frame conjecture is not
used.  This completes the invariant for precisely the marked weighted
OS45 member in the statement.
\end{enumerate}
\end{proof}
